\documentclass[11pt]{article}
% The preamble every project's paper shares: % The preamble every project's paper shares: % The preamble every project's paper shares: \input{../../preamble} from
% projects/<project>/paper/. Boards come from chessboard and xskak, plots from
% pgfplots; every number a paper prints comes from its generated.tex.
\usepackage[LSBC4,T1]{fontenc}
\usepackage[utf8]{inputenc}
\usepackage{lmodern}
\usepackage[margin=1in]{geometry}
\usepackage{microtype}
\usepackage{amsmath}
\usepackage{booktabs}
\usepackage{tabularx}
\usepackage{array}
\usepackage{xcolor}
\usepackage[obeyspaces,spaces]{url}
\usepackage{xskak}
\usepackage{pgfplots}
\pgfplotsset{compat=1.18}
\usepackage{caption}
\usepackage{subcaption}
\usepackage{enumitem}
\usepackage[section]{placeins}
\usepackage[hidelinks]{hyperref}

\setchessboard{boardfontencoding=LSBC4, boardfontfamily=skaknew, blackfieldcolor=black!22,
  whitefieldcolor=white, setfontcolors, showmover=false, label=true, labelfontsize=6pt,
  boardfontsize=14pt}

% A move or line in UCI, and a position in FEN, both breakable at a slash.
\newcommand{\uci}[1]{\texttt{#1}}
\newcommand{\fen}[1]{{\footnotesize\path{#1}}}
\newcolumntype{L}{>{\raggedright\arraybackslash}X}

\setlength{\parskip}{0.4em}
\setlength{\parindent}{0pt}
 from
% projects/<project>/paper/. Boards come from chessboard and xskak, plots from
% pgfplots; every number a paper prints comes from its generated.tex.
\usepackage[LSBC4,T1]{fontenc}
\usepackage[utf8]{inputenc}
\usepackage{lmodern}
\usepackage[margin=1in]{geometry}
\usepackage{microtype}
\usepackage{amsmath}
\usepackage{booktabs}
\usepackage{tabularx}
\usepackage{array}
\usepackage{xcolor}
\usepackage[obeyspaces,spaces]{url}
\usepackage{xskak}
\usepackage{pgfplots}
\pgfplotsset{compat=1.18}
\usepackage{caption}
\usepackage{subcaption}
\usepackage{enumitem}
\usepackage[section]{placeins}
\usepackage[hidelinks]{hyperref}

\setchessboard{boardfontencoding=LSBC4, boardfontfamily=skaknew, blackfieldcolor=black!22,
  whitefieldcolor=white, setfontcolors, showmover=false, label=true, labelfontsize=6pt,
  boardfontsize=14pt}

% A move or line in UCI, and a position in FEN, both breakable at a slash.
\newcommand{\uci}[1]{\texttt{#1}}
\newcommand{\fen}[1]{{\footnotesize\path{#1}}}
\newcolumntype{L}{>{\raggedright\arraybackslash}X}

\setlength{\parskip}{0.4em}
\setlength{\parindent}{0pt}
 from
% projects/<project>/paper/. Boards come from chessboard and xskak, plots from
% pgfplots; every number a paper prints comes from its generated.tex.
\usepackage[LSBC4,T1]{fontenc}
\usepackage[utf8]{inputenc}
\usepackage{lmodern}
\usepackage[margin=1in]{geometry}
\usepackage{microtype}
\usepackage{amsmath}
\usepackage{booktabs}
\usepackage{tabularx}
\usepackage{array}
\usepackage{xcolor}
\usepackage[obeyspaces,spaces]{url}
\usepackage{xskak}
\usepackage{pgfplots}
\pgfplotsset{compat=1.18}
\usepackage{caption}
\usepackage{subcaption}
\usepackage{enumitem}
\usepackage[section]{placeins}
\usepackage[hidelinks]{hyperref}

\setchessboard{boardfontencoding=LSBC4, boardfontfamily=skaknew, blackfieldcolor=black!22,
  whitefieldcolor=white, setfontcolors, showmover=false, label=true, labelfontsize=6pt,
  boardfontsize=14pt}

% A move or line in UCI, and a position in FEN, both breakable at a slash.
\newcommand{\uci}[1]{\texttt{#1}}
\newcommand{\fen}[1]{{\footnotesize\path{#1}}}
\newcolumntype{L}{>{\raggedright\arraybackslash}X}

\setlength{\parskip}{0.4em}
\setlength{\parindent}{0pt}

% Written by projects/band-above/paper/generate.mjs; do not edit by hand.

\newcommand{\AppCommit}{d8b771238c}
\newcommand{\Month}{2026-09}
\newcommand{\HisLo}{1250}
\newcommand{\HisHi}{1449}
\newcommand{\AboveLo}{1450}
\newcommand{\AboveHi}{1649}
\newcommand{\Draws}{2,000}
\newcommand{\Seed}{20261014}
\newcommand{\Discount}{0.6}
\newcommand{\EloPerPoint}{7}
\newcommand{\Causes}{ten}
\newcommand{\PartitionCauses}{seven}
\newcommand{\GamesBuilt}{74,000}
\newcommand{\GamesRead}{74,000}
\newcommand{\BlackChecked}{800}
\newcommand{\PlayerGames}{148,000}
\newcommand{\YardstickK}{2.22}
\newcommand{\LichessK}{3.68}
\newcommand{\YardstickPositions}{2,276,949}
\newcommand{\YardstickPlayerGames}{39,903}
\newcommand{\HisPlayerGames}{39,903}
\newcommand{\AbovePlayerGames}{34,973}
\newcommand{\HisGames}{24,168}
\newcommand{\AboveGames}{20,589}
\newcommand{\BandLo}{850--1049}
\newcommand{\BandHi}{1850--2049}
\newcommand{\Bands}{six}
\newcommand{\LowK}{1.76}
\newcommand{\HighK}{3.04}
\newcommand{\BandTable}{%
\begin{tabular}{lrrrr}
\toprule
player's rating & player-games & games & slope $k$ ($\times 10^{-3}$) & plies \\
\midrule
850--1049 & 13,282 & 8,310 & 1.76 & 60 \\
1050--1249 & 21,692 & 13,937 & 2.01 & 62 \\
1250--1449 (his) & 39,903 & 24,168 & 2.22 & 64 \\
1450--1649 (above) & 34,973 & 20,589 & 2.45 & 66 \\
1650--1849 & 15,268 & 9,353 & 2.72 & 69 \\
1850--2049 & 12,424 & 8,025 & 3.04 & 73 \\
\bottomrule
\end{tabular}
}
\newcommand{\SeparatesTable}{%
\begin{tabular}{lrrllcl}
\toprule
 & \multicolumn{2}{c}{points a game} & & Holm & & share of \\
\cmidrule(lr){2-3}
cause & his band & above & $\Delta$ & $p$ & separates & the difference \\
\midrule
opening, tense & 13.8 & 11.3 & $-$2.51 {\footnotesize ($-$2.79, $-$2.24)} & 0.010 & yes & 45\% {\footnotesize (37, 58)} \\
opening, quiet & 12.1 & 11.3 & $-$0.72 {\footnotesize ($-$0.89, $-$0.53)} & 0.010 & yes & 13\% {\footnotesize (10, 18)} \\
middlegame, tense & 38.2 & 35.1 & $-$3.05 {\footnotesize ($-$3.67, $-$2.42)} & 0.010 & yes & 55\% {\footnotesize (45, 67)} \\
middlegame, quiet & 14.8 & 15.7 & $+$0.97 {\footnotesize ($+$0.66, $+$1.26)} & 0.010 & no & $-$18\% {\footnotesize ($-$26, $-$11)} \\
endgame, tense & 18.9 & 17.9 & $-$1.05 {\footnotesize ($-$1.60, $-$0.48)} & 0.010 & yes & 19\% {\footnotesize (10, 26)} \\
endgame, quiet & 8.8 & 9.4 & $+$0.56 {\footnotesize ($+$0.18, $+$0.93)} & 0.010 & no & $-$10\% {\footnotesize ($-$20, $-$3)} \\
under a minute left & 2.2 & 2.4 & $+$0.24 {\footnotesize ($+$0.02, $+$0.46)} & 0.033 & no & $-$4\% {\footnotesize ($-$9, $-$0)} \\
\midrule
tense, under 5 seconds & 23.8 & 22.4 & $-$1.40 {\footnotesize ($-$1.83, $-$0.97)} & 0.010 & yes &  \\
after the opponent's error & 28.5 & 25.5 & $-$2.97 {\footnotesize ($-$3.56, $-$2.35)} & 0.010 & yes &  \\
from move 40 & 8.9 & 9.8 & $+$0.85 {\footnotesize ($+$0.31, $+$1.39)} & 0.010 & no &  \\
\midrule
every move & 108.6 & 103.1 & $-$5.55 {\footnotesize ($-$6.70, $-$4.41)} &  &  &  \\
\bottomrule
\end{tabular}
}
\newcommand{\Separate}{six}
\newcommand{\SeparateCap}{Six}
\newcommand{\RankedList}{tense middlegames, the move after the opponent's serious error, tense openings, tense positions answered in under five seconds, tense endgames and quiet openings}
\newcommand{\AllGap}{5.55}
\newcommand{\AllGapLo}{4.41}
\newcommand{\AllGapHi}{6.70}
\newcommand{\OpeningTenseHis}{13.8}
\newcommand{\OpeningTenseAbove}{11.3}
\newcommand{\OpeningTenseDelta}{$-$2.51}
\newcommand{\OpeningTenseDeltaLo}{$-$2.79}
\newcommand{\OpeningTenseDeltaHi}{$-$2.24}
\newcommand{\OpeningTenseShare}{45}
\newcommand{\OpeningQuietHis}{12.1}
\newcommand{\OpeningQuietAbove}{11.3}
\newcommand{\OpeningQuietDelta}{$-$0.72}
\newcommand{\OpeningQuietDeltaLo}{$-$0.89}
\newcommand{\OpeningQuietDeltaHi}{$-$0.53}
\newcommand{\OpeningQuietShare}{13}
\newcommand{\MiddlegameTenseHis}{38.2}
\newcommand{\MiddlegameTenseAbove}{35.1}
\newcommand{\MiddlegameTenseDelta}{$-$3.05}
\newcommand{\MiddlegameTenseDeltaLo}{$-$3.67}
\newcommand{\MiddlegameTenseDeltaHi}{$-$2.42}
\newcommand{\MiddlegameTenseShare}{55}
\newcommand{\MiddlegameQuietHis}{14.8}
\newcommand{\MiddlegameQuietAbove}{15.7}
\newcommand{\MiddlegameQuietDelta}{$+$0.97}
\newcommand{\MiddlegameQuietDeltaLo}{$+$0.66}
\newcommand{\MiddlegameQuietDeltaHi}{$+$1.26}
\newcommand{\MiddlegameQuietShare}{$-$18}
\newcommand{\EndgameTenseHis}{18.9}
\newcommand{\EndgameTenseAbove}{17.9}
\newcommand{\EndgameTenseDelta}{$-$1.05}
\newcommand{\EndgameTenseDeltaLo}{$-$1.60}
\newcommand{\EndgameTenseDeltaHi}{$-$0.48}
\newcommand{\EndgameTenseShare}{19}
\newcommand{\EndgameQuietHis}{8.8}
\newcommand{\EndgameQuietAbove}{9.4}
\newcommand{\EndgameQuietDelta}{$+$0.56}
\newcommand{\EndgameQuietDeltaLo}{$+$0.18}
\newcommand{\EndgameQuietDeltaHi}{$+$0.93}
\newcommand{\EndgameQuietShare}{$-$10}
\newcommand{\ClockLowHis}{2.2}
\newcommand{\ClockLowAbove}{2.4}
\newcommand{\ClockLowDelta}{$+$0.24}
\newcommand{\ClockLowDeltaLo}{$+$0.02}
\newcommand{\ClockLowDeltaHi}{$+$0.46}
\newcommand{\ClockLowShare}{$-$4}
\newcommand{\FastTenseHis}{23.8}
\newcommand{\FastTenseAbove}{22.4}
\newcommand{\FastTenseDelta}{$-$1.40}
\newcommand{\FastTenseDeltaLo}{$-$1.83}
\newcommand{\FastTenseDeltaHi}{$-$0.97}
\newcommand{\AfterGiftHis}{28.5}
\newcommand{\AfterGiftAbove}{25.5}
\newcommand{\AfterGiftDelta}{$-$2.97}
\newcommand{\AfterGiftDeltaLo}{$-$3.56}
\newcommand{\AfterGiftDeltaHi}{$-$2.35}
\newcommand{\LateHis}{8.9}
\newcommand{\LateAbove}{9.8}
\newcommand{\LateDelta}{$+$0.85}
\newcommand{\LateDeltaLo}{$+$0.31}
\newcommand{\LateDeltaHi}{$+$1.39}
\newcommand{\AllHis}{108.6}
\newcommand{\AllAbove}{103.1}
\newcommand{\AllDelta}{$-$5.55}
\newcommand{\AllDeltaLo}{$-$6.70}
\newcommand{\AllDeltaHi}{$-$4.41}
\newcommand{\TenseShare}{119}
\newcommand{\GiveBack}{32}
\newcommand{\PerMoveTable}{%
\begin{tabular}{lrrlll}
\toprule
 & \multicolumn{2}{c}{moves a game} & \multicolumn{3}{c}{points a move} \\
\cmidrule(lr){2-3}\cmidrule(lr){4-6}
cause & his band & above & his band & above & $\Delta$ \\
\midrule
opening, tense & 4.5 & 4.2 & 3.07 {\footnotesize (3.03, 3.10)} & 2.66 {\footnotesize (2.62, 2.69)} & $-$0.41 {\footnotesize ($-$0.46, $-$0.36)} \\
opening, quiet & 5.8 & 6.1 & 2.07 {\footnotesize (2.05, 2.09)} & 1.86 {\footnotesize (1.84, 1.88)} & $-$0.21 {\footnotesize ($-$0.24, $-$0.19)} \\
middlegame, tense & 8.5 & 8.5 & 4.47 {\footnotesize (4.42, 4.51)} & 4.13 {\footnotesize (4.09, 4.17)} & $-$0.33 {\footnotesize ($-$0.39, $-$0.28)} \\
middlegame, quiet & 3.5 & 4.1 & 4.17 {\footnotesize (4.12, 4.21)} & 3.87 {\footnotesize (3.83, 3.91)} & $-$0.29 {\footnotesize ($-$0.36, $-$0.24)} \\
endgame, tense & 5.6 & 5.7 & 3.35 {\footnotesize (3.30, 3.39)} & 3.16 {\footnotesize (3.11, 3.21)} & $-$0.19 {\footnotesize ($-$0.26, $-$0.12)} \\
endgame, quiet & 2.9 & 3.1 & 3.03 {\footnotesize (2.97, 3.09)} & 2.99 {\footnotesize (2.93, 3.05)} & $-$0.04 {\footnotesize ($-$0.12, $+$0.04)} \\
under a minute left & 0.6 & 0.7 & 3.71 {\footnotesize (3.53, 3.89)} & 3.49 {\footnotesize (3.31, 3.67)} & $-$0.22 {\footnotesize ($-$0.46, $+$0.02)} \\
\midrule
tense, under 5 seconds & 8.9 & 9.1 & 2.67 {\footnotesize (2.64, 2.69)} & 2.45 {\footnotesize (2.42, 2.47)} & $-$0.22 {\footnotesize ($-$0.26, $-$0.18)} \\
after the opponent's error & 3.6 & 3.5 & 7.88 {\footnotesize (7.80, 7.96)} & 7.35 {\footnotesize (7.27, 7.43)} & $-$0.53 {\footnotesize ($-$0.64, $-$0.42)} \\
from move 40 & 3.1 & 3.4 & 2.90 {\footnotesize (2.82, 2.97)} & 2.88 {\footnotesize (2.81, 2.96)} & $-$0.02 {\footnotesize ($-$0.12, $+$0.09)} \\
\midrule
every move & 31.5 & 32.4 & 3.44 {\footnotesize (3.42, 3.47)} & 3.18 {\footnotesize (3.16, 3.20)} & $-$0.26 {\footnotesize ($-$0.29, $-$0.23)} \\
\bottomrule
\end{tabular}
}
\newcommand{\PerMoveLess}{seven}
\newcommand{\HisPlies}{64.4}
\newcommand{\AbovePlies}{66.1}
\newcommand{\HisQuietMiddlegameMoves}{3.5}
\newcommand{\AboveQuietMiddlegameMoves}{4.1}
\newcommand{\AllPerMoveHis}{3.44}
\newcommand{\AllPerMoveAbove}{3.18}
\newcommand{\GiftPerMoveHis}{7.9}
\newcommand{\GiftPerMoveAbove}{7.3}
\newcommand{\GiftRatio}{2.3}
\newcommand{\DiscountTable}{%
\begin{tabular}{ll}
\toprule
cause & $\delta$ \\
\midrule
opening, tense & 0.30 {\footnotesize (0.28, 0.32)} \\
opening, quiet & 0.29 {\footnotesize (0.25, 0.32)} \\
middlegame, tense & 0.29 {\footnotesize (0.28, 0.30)} \\
middlegame, quiet & 0.24 {\footnotesize (0.22, 0.26)} \\
endgame, tense & 0.22 {\footnotesize (0.21, 0.23)} \\
endgame, quiet & 0.16 {\footnotesize (0.15, 0.18)} \\
under a minute left & 0.31 {\footnotesize (0.27, 0.35)} \\
\midrule
tense, under 5 seconds & 0.33 {\footnotesize (0.32, 0.35)} \\
after the opponent's error & 0.10 {\footnotesize (0.09, 0.11)} \\
from move 40 & 0.12 {\footnotesize (0.11, 0.13)} \\
\midrule
every move & 0.22 {\footnotesize (0.21, 0.23)} \\
\bottomrule
\end{tabular}
}
\newcommand{\DiscountMin}{0.10}
\newcommand{\DiscountMax}{0.33}
\newcommand{\AllDiscount}{0.22}
\newcommand{\AllDiscountLo}{0.21}
\newcommand{\AllDiscountHi}{0.23}
\newcommand{\TheirsTable}{%
\begin{tabular}{lrll}
\toprule
 & correlation of & \multicolumn{2}{c}{$\delta$ with the opponent's points} \\
\cmidrule(lr){3-4}
cause & the two players' points & the cause alone & every partition cause \\
\midrule
opening, tense & 0.61 & 0.78 {\footnotesize (0.75, 0.82)} & 1.06 {\footnotesize (1.04, 1.07)} \\
opening, quiet & 0.52 & 0.61 {\footnotesize (0.55, 0.67)} & 1.03 {\footnotesize (1.01, 1.06)} \\
middlegame, tense & 0.66 & 0.85 {\footnotesize (0.83, 0.87)} & 1.06 {\footnotesize (1.06, 1.07)} \\
middlegame, quiet & 0.56 & 0.53 {\footnotesize (0.49, 0.57)} & 1.07 {\footnotesize (1.05, 1.08)} \\
endgame, tense & 0.69 & 0.71 {\footnotesize (0.69, 0.74)} & 1.02 {\footnotesize (1.01, 1.03)} \\
endgame, quiet & 0.69 & 0.54 {\footnotesize (0.50, 0.58)} & 1.02 {\footnotesize (1.01, 1.04)} \\
under a minute left & 0.24 & 0.41 {\footnotesize (0.38, 0.44)} & 1.25 {\footnotesize (1.22, 1.27)} \\
\midrule
tense, under 5 seconds & 0.39 & 0.55 {\footnotesize (0.53, 0.57)} &  \\
after the opponent's error & 0.67 & 0.30 {\footnotesize (0.28, 0.32)} &  \\
from move 40 & 0.87 & 0.95 {\footnotesize (0.92, 0.98)} &  \\
\midrule
every move & 0.79 & 1.05 {\footnotesize (1.05, 1.06)} &  \\
\bottomrule
\end{tabular}
}
\newcommand{\CorrAll}{0.79}
\newcommand{\TheirsAll}{1.05}
\newcommand{\TheirsAllLo}{1.05}
\newcommand{\TheirsAllHi}{1.06}
\newcommand{\TheirsAllOpp}{1.05}
\newcommand{\PartitionMin}{1.02}
\newcommand{\PartitionMax}{1.25}
\newcommand{\GiftTheirs}{0.30}
\newcommand{\KindsTable}{%
\setlength{\tabcolsep}{4pt}
\begin{tabular}{lllllr}
\toprule
 & \multicolumn{4}{c}{per game} & hanging a piece: \\
\cmidrule(lr){2-5}
player's rating & serious errors & hanging a piece & mate in one & a forced mate & once every $N$ games \\
\midrule
850--1049 & 4.07 {\footnotesize (4.00, 4.13)} & 1.02 {\footnotesize (1.00, 1.04)} & 0.072 {\footnotesize (0.067, 0.077)} & 0.072 {\footnotesize (0.067, 0.077)} & 1.0 \\
1050--1249 & 3.81 {\footnotesize (3.76, 3.86)} & 0.90 {\footnotesize (0.89, 0.92)} & 0.060 {\footnotesize (0.057, 0.063)} & 0.066 {\footnotesize (0.062, 0.069)} & 1.1 \\
1250--1449 & 3.67 {\footnotesize (3.63, 3.70)} & 0.82 {\footnotesize (0.81, 0.83)} & 0.048 {\footnotesize (0.046, 0.050)} & 0.062 {\footnotesize (0.060, 0.065)} & 1.2 \\
1450--1649 & 3.51 {\footnotesize (3.48, 3.55)} & 0.74 {\footnotesize (0.73, 0.75)} & 0.042 {\footnotesize (0.040, 0.045)} & 0.055 {\footnotesize (0.053, 0.058)} & 1.4 \\
1650--1849 & 3.39 {\footnotesize (3.34, 3.44)} & 0.66 {\footnotesize (0.65, 0.68)} & 0.031 {\footnotesize (0.028, 0.034)} & 0.048 {\footnotesize (0.044, 0.052)} & 1.5 \\
1850--2049 & 3.16 {\footnotesize (3.11, 3.21)} & 0.59 {\footnotesize (0.58, 0.61)} & 0.023 {\footnotesize (0.020, 0.026)} & 0.044 {\footnotesize (0.040, 0.048)} & 1.7 \\
\bottomrule
\end{tabular}
}
\newcommand{\KindDiffTable}{%
\begin{tabular}{lrrll}
\toprule
 & \multicolumn{2}{c}{points a game} & & share of \\
\cmidrule(lr){2-3}
kind & his band & above & $\Delta$ & the difference \\
\midrule
serious errors & 66.6 & 61.4 & $-$5.20 {\footnotesize ($-$6.18, $-$4.23)} & 94\% {\footnotesize (87, 101)} \\
hanging a piece & 16.6 & 14.2 & $-$2.38 {\footnotesize ($-$2.70, $-$2.07)} & 43\% {\footnotesize (36, 52)} \\
allowing mate in one & 2.4 & 2.0 & $-$0.38 {\footnotesize ($-$0.54, $-$0.21)} & 7\% {\footnotesize (4, 10)} \\
allowing a forced mate & 2.8 & 2.4 & $-$0.35 {\footnotesize ($-$0.54, $-$0.18)} & 6\% {\footnotesize (3, 10)} \\
other serious errors & 44.9 & 42.8 & $-$2.09 {\footnotesize ($-$2.86, $-$1.34)} & 38\% {\footnotesize (29, 45)} \\
\bottomrule
\end{tabular}
}
\newcommand{\HisSerious}{3.67}
\newcommand{\HisHanging}{0.82}
\newcommand{\HisHangingEvery}{1.2}
\newcommand{\HisMateOneEvery}{20.8}
\newcommand{\HisMateNEvery}{16.1}
\newcommand{\AboveSerious}{3.51}
\newcommand{\AboveHanging}{0.74}
\newcommand{\AboveHangingEvery}{1.4}
\newcommand{\AboveMateOneEvery}{23.6}
\newcommand{\AboveMateNEvery}{18.1}
\newcommand{\ThinkGiftErrors}{33}
\newcommand{\ThinkGiftMoves}{14}
\newcommand{\SeriousDiffShare}{94}
\newcommand{\HangingDiffShare}{43}
\newcommand{\MateOneDiffShare}{7}
\newcommand{\MateNDiffShare}{6}
\newcommand{\UntaggedDiffShare}{38}
\newcommand{\MatesDiffShare}{13}
\newcommand{\TrendTable}{%
\begin{tabular}{lrrrrrr}
\toprule
cause & 850--1049 & 1050--1249 & 1250--1449 & 1450--1649 & 1650--1849 & 1850--2049 \\
\midrule
opening, tense & 24.1 & 17.6 & 13.8 & 11.3 & 8.5 & 6.8 \\
opening, quiet & 14.1 & 13.2 & 12.1 & 11.3 & 10.3 & 9.2 \\
middlegame, tense & 49.6 & 41.9 & 38.2 & 35.1 & 32.6 & 29.0 \\
middlegame, quiet & 12.8 & 13.6 & 14.8 & 15.7 & 16.7 & 16.7 \\
endgame, tense & 18.6 & 19.0 & 18.9 & 17.9 & 17.5 & 16.5 \\
endgame, quiet & 6.4 & 7.7 & 8.8 & 9.4 & 10.0 & 10.4 \\
under a minute left & 2.0 & 2.2 & 2.2 & 2.4 & 2.9 & 4.1 \\
\midrule
tense, under 5 seconds & 27.8 & 25.0 & 23.8 & 22.4 & 20.9 & 19.0 \\
after the opponent's error & 39.7 & 32.4 & 28.5 & 25.5 & 22.7 & 19.0 \\
from move 40 & 5.7 & 7.4 & 8.9 & 9.8 & 11.4 & 13.1 \\
\midrule
every move & 127.5 & 115.3 & 108.6 & 103.1 & 98.4 & 92.8 \\
\bottomrule
\end{tabular}
}
\newcommand{\AllLow}{127.5}
\newcommand{\AllHigh}{92.8}
\newcommand{\PliesLow}{60}
\newcommand{\PliesHigh}{73}
\newcommand{\AllPlot}{%
\begin{tikzpicture}
\begin{axis}[scale only axis, width=0.75\linewidth, height=3.6cm, xtick={0,...,5}, xticklabels={950,1150,1350,1550,1750,1950},
  xmin=-0.5, xmax=5.5, x tick label style={font=\scriptsize}, y tick label style={font=\scriptsize}, ylabel={points a game}, xlabel={player's rating, the band's centre}, grid=major]
\addplot[thick, black, mark=*, mark size=1.5pt, error bars/.cd, y dir=both, y explicit] coordinates {
(0, 127.499) += (0, 1.729) -= (0, 1.757)
(1, 115.324) += (0, 1.222) -= (0, 1.174)
(2, 108.640) += (0, 0.822) -= (0, 0.843)
(3, 103.087) += (0, 0.872) -= (0, 0.844)
(4, 98.419) += (0, 1.195) -= (0, 1.128)
(5, 92.778) += (0, 1.220) -= (0, 1.183)
};
\end{axis}
\end{tikzpicture}
}
\newcommand{\SeriousPlot}{%
\begin{tikzpicture}
\begin{axis}[scale only axis, width=0.75\linewidth, height=3.6cm, xtick={0,...,5}, xticklabels={950,1150,1350,1550,1750,1950},
  xmin=-0.5, xmax=5.5, x tick label style={font=\scriptsize}, y tick label style={font=\scriptsize}, ylabel={serious errors a game}, xlabel={player's rating, the band's centre}, grid=major]
\addplot[thick, black, mark=*, mark size=1.5pt, error bars/.cd, y dir=both, y explicit] coordinates {
(0, 4.067) += (0, 0.061) -= (0, 0.064)
(1, 3.810) += (0, 0.046) -= (0, 0.045)
(2, 3.666) += (0, 0.032) -= (0, 0.035)
(3, 3.515) += (0, 0.036) -= (0, 0.036)
(4, 3.389) += (0, 0.049) -= (0, 0.048)
(5, 3.163) += (0, 0.051) -= (0, 0.049)
};
\end{axis}
\end{tikzpicture}
}
\newcommand{\ArchiveGames}{633}
\newcommand{\ArchivePlies}{63}
\newcommand{\SplitTable}{%
\begin{tabular}{llll}
\toprule
cause & his games & his band & the band above \\
\midrule
opening, tense & 9.6\% {\footnotesize (8.5, 10.8)} & 12.7\% {\footnotesize (12.5, 12.9)} & 10.9\% {\footnotesize (10.7, 11.1)} \\
opening, quiet & 10.3\% {\footnotesize (9.4, 11.1)} & 11.1\% {\footnotesize (11.0, 11.2)} & 11.0\% {\footnotesize (10.8, 11.1)} \\
middlegame, tense & 36.3\% {\footnotesize (34.2, 38.6)} & 35.1\% {\footnotesize (34.8, 35.5)} & 34.1\% {\footnotesize (33.7, 34.4)} \\
middlegame, quiet & 13.2\% {\footnotesize (12.0, 14.3)} & 13.6\% {\footnotesize (13.4, 13.8)} & 15.3\% {\footnotesize (15.1, 15.5)} \\
endgame, tense & 20.6\% {\footnotesize (18.7, 22.5)} & 17.4\% {\footnotesize (17.1, 17.7)} & 17.3\% {\footnotesize (17.0, 17.6)} \\
endgame, quiet & 7.8\% {\footnotesize (6.6, 9.1)} & 8.1\% {\footnotesize (7.9, 8.3)} & 9.1\% {\footnotesize (8.8, 9.3)} \\
under a minute left & 2.3\% {\footnotesize (1.4, 3.3)} & 2.0\% {\footnotesize (1.9, 2.1)} & 2.4\% {\footnotesize (2.2, 2.5)} \\
\bottomrule
\end{tabular}
}
\newcommand{\HisKindsTable}{%
\begin{tabular}{llllr}
\toprule
 & \multicolumn{3}{c}{per game} & his games: \\
\cmidrule(lr){2-4}
kind & his games & his band & the band above & once every $N$ \\
\midrule
hanging a piece & 0.777 {\footnotesize (0.707, 0.849)} & 0.823 {\footnotesize (0.814, 0.833)} & 0.738 {\footnotesize (0.728, 0.748)} & 1.3 \\
allowing mate in one & 0.073 {\footnotesize (0.053, 0.093)} & 0.048 {\footnotesize (0.046, 0.050)} & 0.042 {\footnotesize (0.040, 0.045)} & 13.8 \\
allowing a forced mate & 0.100 {\footnotesize (0.073, 0.127)} & 0.062 {\footnotesize (0.060, 0.065)} & 0.055 {\footnotesize (0.053, 0.058)} & 10.0 \\
\bottomrule
\end{tabular}
}
\newcommand{\SplitEndgameTenseMine}{20.6}
\newcommand{\SplitEndgameTenseHis}{17.4}
\newcommand{\SplitEndgameTenseAbove}{17.3}
\newcommand{\SplitOpeningTenseMine}{9.6}
\newcommand{\SplitOpeningTenseHis}{12.7}
\newcommand{\SplitOpeningTenseAbove}{10.9}
\newcommand{\ArchiveHangingEvery}{1.3}
\newcommand{\ArchiveMateOneEvery}{13.8}
\newcommand{\ArchiveMateNEvery}{10.0}
\newcommand{\MateRatio}{1.6}


\title{What does the band above do differently?\\[0.3em]
\large The leak ledger's causes on \GamesBuilt{} Lichess 10-minute games, his level against 200 points above}
\author{Chess Coach research lab}
\date{}

\begin{document}
\maketitle

\begin{abstract}
Players 200 points above his level give away \AllGap{} fewer points a game than players at his
level (\AllGapLo{} to \AllGapHi), on \GamesBuilt{} Lichess 10-minute games built with the app's
own code. \SeparateCap{} of the leak ledger's \Causes{} causes separate the two bands. Ranked by the
difference, they are \RankedList. Tense positions carry more than the whole difference
(\TenseShare\%). The band above gives away more a game in quiet middlegames and endgames only
because its games are longer: per move it gives away less or the same on every cause. Serious
errors carry \SeriousDiffShare\% of the difference, and hanging a piece \HangingDiffShare\%; his
level hangs one once every \HisHangingEvery{} games, the band above once every
\AboveHangingEvery. The rule's fit of the result on the points given away finds the app's
\Discount{} discount too high on every cause. But that fit leaves out the opponent's points, which
rise and fall with the player's. With them in, in a check added after the results were read, a
point given away costs about a point of result (\TheirsAll). The data give no reason to lower the
\Discount. In his own games, tense endgames take a larger share of his points than at either level,
and he allows a mate about \MateRatio{} times as often as players at his level.
\end{abstract}

\section{The question}

The app's ONE focus is a leak from the metrics service's ledger (\texttt{leakLedgerOf}, at commit
\texttt{\AppCommit}): the points a game the player gives away on one kind of move, against what his
opponents give away on the same kind of move in the same games. The kinds of move are the ledger's
\Causes{} causes: the phase and whether the position was tense, under a minute on the clock, a
tense position answered in under five seconds, the move after the opponent's serious error, and
moves from move 40. A point is a percent of a game, the drop in expected score across the move,
and a leak is priced in Elo as \EloPerPoint{} times a hand-set discount of \Discount{} times the
gap: the share of a point given away that a fix would bank. Two things the ledger does not
measure: where players a band above his give away less than players at his level, and how much of
a point given away shows in the result. So:
\begin{enumerate}[nosep]
\item On which causes do players 200 points above his level (\AboveLo--\AboveHi{} on Lichess) give
away fewer points a game than players at his level (\HisLo--\HisHi), and which carry the difference?
\item How much does a point given away on each cause cost in the result, against the app's
\Discount?
\item How often do players at each level hang a piece or allow a mate, and what does it cost them?
\end{enumerate}

\section{The data}

The games are the \GamesBuilt{} rated 10-minute Lichess games the think-time study kept from
\Month{} (CC0), with Lichess's evaluation on every move and the clock. Each was built once with the
app's own code through the lab's bridge: the metrics service's \texttt{primaryFacts} with the
routine's ports, and the review service's reader for its motif tags. Each game gives two
player-games, one for each player as ``me''; Black's facts are derived from White's build with the
evaluations and the result turned around, and on \BlackChecked{} games they matched a real Black
build, move facts, score and tags. A player-game's band is the player's Lichess rating
(Table~\ref{tab:bands}). His chess.com rapid rating comes to the middle of his band by the app's
default conversion.

\begin{table}[h]
\centering
\small\BandTable
\caption{The bands. The slope is the band curve's, fitted on the band's own player-games; plies
are a game's mean length.}
\label{tab:bands}
\end{table}

\section{The method}

Every measure is the app's. The points a player gives away on a cause are the method service's
\texttt{causeLoss} for that cause: the drop in expected score across each of his moves of that
kind, on the band curve, floored at zero. One yardstick prices every band: the curve's slope fitted
by the metrics service's \texttt{fitWinCurve} on his band's player-games ($k = \YardstickK \times
10^{-3}$ per centipawn, from \YardstickPositions{} positions; the Lichess curve's is \LichessK),
so the same move costs the same points at every level. The ledger's first seven causes, the six
phase-by-position ones and the low clock, split every move between them, so their differences add
up to the difference on every move.

Error kinds come from the review reader's motif tags on the moves it calls a mistake or a blunder,
read from the evaluations alone: with no engine line, a move that leaves one of the mover's pieces
newly hanging is \emph{hanging a piece}, and the evaluation or the board gives \emph{allowing mate in
one} and \emph{allowing a forced mate} (in two to five). The three kinds and the untagged rest split
the serious errors.

Means are over player-games. Intervals come from \Draws{} bootstrap draws over games (a game's two
player-games share its Poisson weight; seed \Seed), the 95\% percentile interval, with a $p$-value
of twice the smaller share of draws on either side of zero.

The rule, committed before any result was computed. (1) A cause \emph{separates the bands} when the
band above's mean points a game on it minus his band's, $\Delta_c$, is below zero with a
Holm-adjusted $p$ under 0.05 over the \Causes{} causes; the causes that separate them, ranked by
$\Delta_c$, are the priorities the data give his band. (2) Over his band's player-games, a least
squares fit of the result (0, 50 or 100) on the points given away on a cause and the rating
difference gives $\delta_c$, minus the slope on the points: the share of a point given away that
shows in the result. The app's \Discount{} is \emph{too high} for a cause when $\delta_c$'s interval
lies below it, \emph{too low} when above, and \emph{within} otherwise.

Four checks were added after the results were read, and each is marked where it appears: points per
move, the error kinds' share of the difference, the discount fitted with the opponent's points, and
his own games' split. They test nothing and carry no rule; they say what the registered numbers mean.

\section{Results}

\subsection{What separates the bands}

Over every move, the band above gives away \AllAbove{} points a game and his band \AllHis, a
difference of \AllDelta{} (\AllDeltaLo{} to \AllDeltaHi; Table~\ref{tab:separates}). \SeparateCap{}
of the \Causes{} causes separate the bands. Ranked by the difference, they are \RankedList. The first
two cost about three points a game each: tense middlegames \MiddlegameTenseDelta{} and the move
after the opponent's serious error \AfterGiftDelta.

The \PartitionCauses{} causes that split every move add up to the difference on every move, and the
tense ones carry more than all of it: \OpeningTenseShare\% in the opening, \MiddlegameTenseShare\%
in the middlegame and \EndgameTenseShare\% in the endgame, \TenseShare\% together. Quiet openings
add \OpeningQuietShare\%. The band above gives away \emph{more} a game in quiet middlegames
(\MiddlegameQuietDelta), quiet endgames (\EndgameQuietDelta), under a minute on the clock
(\ClockLowDelta) and from move 40 (\LateDelta), the four causes that do not separate the bands; the
first three give back \GiveBack\% of the difference.

\begin{table}[h]
\centering
\scriptsize\SeparatesTable
\caption{What separates the bands: points a game given away on each cause at his level
(\HisLo--\HisHi) and the band above (\AboveLo--\AboveHi), the difference $\Delta$ (above minus his)
with its 95\% interval, the Holm-adjusted $p$, the rule's verdict, and each partition cause's share
of the difference on every move. The first seven causes split every move; the next three overlap
them.}
\label{tab:separates}
\end{table}

\subsection{Per game and per move}

Why would stronger players give away more in quiet positions? A cause's points a game are its points
a move times its moves a game, and the band above's games are longer (\AbovePlies{} plies against
\HisPlies), with more quiet middlegame moves in them (\AboveQuietMiddlegameMoves{} a game against
\HisQuietMiddlegameMoves). Per move, a check added after the results were read, the band above gives
away less on \PerMoveLess{} of the \Causes{} causes and the same, within the interval, on quiet
endgames, under a minute on the clock and from move 40 (Table~\ref{tab:permove}); over every move,
\AllPerMoveAbove{} points a move against \AllPerMoveHis. Where the band above gives away more a game,
it meets more of those moves; it does not play them worse. The move after the opponent's serious
error costs the most a move in both bands: \GiftPerMoveHis{} points a move at his level, \GiftRatio{}
times an average move.

\begin{table}[h]
\centering
\scriptsize\PerMoveTable
\caption{Per game and per move (added after the results were read; it tests nothing): each cause's
moves a game, and the points a move given away on it, the band's total points over its total moves,
with the difference (above minus his).}
\label{tab:permove}
\end{table}

\subsection{Error kinds}

Every level gives away fewer points than the one below it, from \AllLow{} points a game at
\BandLo{} to \AllHigh{} at \BandHi, and makes fewer serious errors (Figure~\ref{fig:trend},
Table~\ref{tab:kinds}), while its games get longer, from \PliesLow{} plies to \PliesHigh. Points a
game fall at every band on tense openings and middlegames, quiet openings, fast tense answers and
the move after the opponent's serious error, and rise at every band on quiet middlegames and
endgames and from move 40 (Table~\ref{tab:trend}).

A player at his level makes \HisSerious{} serious errors a game and one at the band above
\AboveSerious. In a check added after the results were read, serious errors carry
\SeriousDiffShare\% of the difference between the two bands and hanging a piece \HangingDiffShare\%
on its own (Table~\ref{tab:kinddiff}): his band hangs a piece once every \HisHangingEvery{} games,
the band above once every \AboveHangingEvery. Allowing a mate is rare at both levels, mate in one
once every \HisMateOneEvery{} games at his and once every \AboveMateOneEvery{} above, and mates carry
\MatesDiffShare\% of the difference. The serious errors the reader leaves untagged, which include
every fork or pin that wins material, carry \UntaggedDiffShare\%.

\begin{figure}[h]
\centering
\begin{subfigure}{0.49\textwidth}\centering\AllPlot\caption{Points a game, every move.}\end{subfigure}\hfill
\begin{subfigure}{0.49\textwidth}\centering\SeriousPlot\caption{Serious errors a game.}\end{subfigure}
\caption{Across the \Bands{} bands, with 95\% intervals (most are narrower than the marks).}
\label{fig:trend}
\end{figure}

\begin{table}[h]
\centering
\scriptsize\KindsTable
\caption{Error kinds by band: serious errors a game and the kinds the reader tags among them, with
95\% intervals.}
\label{tab:kinds}
\end{table}

\begin{table}[h]
\centering
\scriptsize\KindDiffTable
\caption{What each kind costs, his band against the band above (added after the results were read;
it tests nothing): points a game, the difference and its share of the difference on every move.
The four kinds below serious errors split them.}
\label{tab:kinddiff}
\end{table}

\begin{table}[h]
\centering
\scriptsize\TrendTable
\caption{Points a game on each cause, by band.}
\label{tab:trend}
\end{table}

\subsection{The discount}
\label{sec:discount}

By the rule, the app's \Discount{} is too high on all \Causes{} causes: $\delta$ runs from
\DiscountMin{} to \DiscountMax, and is \AllDiscount{} over every move (Table~\ref{tab:discount}).
The fit cannot carry that verdict. On a calibrated curve a game's result is, near enough, 50 plus
the points the opponent gives away less the points the player gives away, and games where one
player gives away more are games where the other does too: over every move, the two players' points
a game correlate at \CorrAll. Left out of the fit, the opponent's points pull the slope on the
player's points towards zero.

With them in, a check added after the results were read and not part of the rule, a point given
away costs \TheirsAll{} points of result over every move (\TheirsAllLo{} to \TheirsAllHi), and a
point the opponent gives away is worth \TheirsAllOpp{} (Table~\ref{tab:theirs}). With the
\PartitionCauses{} partition causes in one fit, each beside the opponent's points on it, a point
given away costs between \PartitionMin{} and \PartitionMax{} points of result on every one. The
other three causes overlap the partition, so a fit of one of them alone still leaves out points that
rise and fall with it, and comes in lower (\GiftTheirs{} for the move after the opponent's serious
error).

So the results give no reason to lower the \Discount: a point given away shows in the result one
for one, as the curve says it should. What the \Discount{} can stand for is the share of a leak a fix
removes, and bands of other players cannot measure that.

\begin{table}[h]
\centering
\small\DiscountTable
\caption{The rule's discount fit, over his band's player-games: $\delta$, minus the slope of the
result on the points given away with the rating difference beside them, and its 95\% interval. By
the rule every cause is \emph{too high} against the app's \Discount.}
\label{tab:discount}
\end{table}

\begin{table}[h]
\centering
\scriptsize\TheirsTable
\caption{The discount with the opponent's points (added after the results were read; not part of the
rule): the correlation of the two players' points a game, $\delta$ with the opponent's points on the
same cause beside the player's, and $\delta$ in one fit with both players' points on every partition
cause.}
\label{tab:theirs}
\end{table}

\subsection{His games}

His \ArchiveGames{} ten-minute games come from the app's own reviews, with its engine and depth, not
Lichess's server analysis, so their points are not on the bands' footing. This section compares only
how they split between the causes and how often a board shows an error; both were added after the
results were read and test nothing.

His points split like his band's on most causes, and two differ (Table~\ref{tab:split}). Tense
endgames take \SplitEndgameTenseMine\% of his points, against \SplitEndgameTenseHis\% at his level
and \SplitEndgameTenseAbove\% above. Tense openings take \SplitOpeningTenseMine\%, against
\SplitOpeningTenseHis\% at his level, about the band above's \SplitOpeningTenseAbove\%. He hangs a
piece about as often as either band, once every \ArchiveHangingEvery{} games, and allows a mate more
often than either: mate in one once every \ArchiveMateOneEvery{} games and a forced mate once every
\ArchiveMateNEvery, about \MateRatio{} times his level's rate (Table~\ref{tab:hiskinds}).

\begin{table}[h]
\centering
\small\SplitTable
\caption{How the points given away split between the partition's causes (added after the results
were read; it tests nothing): his games, his band and the band above, with 95\% intervals.}
\label{tab:split}
\end{table}

\begin{table}[h]
\centering
\scriptsize\HisKindsTable
\caption{How often a board shows an error (added after the results were read; it tests nothing):
per game in his games and in each band, with 95\% intervals.}
\label{tab:hiskinds}
\end{table}

\section{What it means for the app}

\begin{enumerate}[leftmargin=*]
\item \textbf{Where his level loses to the band above.} The ledger's ONE focus picks the cause where
he gives away most against his own opponents. The bands add a second yardstick, where players 200
points above give away less, and for his level it puts tense positions first: the middlegame, the
opening, then the endgame, with the move after the opponent's serious error and fast answers in
tense positions beside them. Quiet middlegames and endgames are not where the band above is better.
\item \textbf{A yardstick from other players' games is per move.} Per game, the band above's longer
games make quiet positions and late moves look like its weakness. The ledger's own comparison,
against the opponent in the same games, meets the same moves and needs no change.
\item \textbf{The move after the opponent's error.} It costs \GiftRatio{} times an average move at
his level and separates the bands, and the think-time study found \ThinkGiftErrors\% of serious
errors there, on \ThinkGiftMoves\% of moves.
The coach can teach a check for that move; the think-time study could not show that more time there
pays.
\item \textbf{Keep the \Discount.} A point given away costs about a point of result. The \Discount{}
can only discount how much of a leak a fix removes, which needs his own games before and after a
focus, not other players' games.
\item \textbf{Rates the coach can cite.} At his level a player hangs a piece once every
\HisHangingEvery{} games and the band above once every \AboveHangingEvery; his own rate can sit
beside them. In his games, mates allowed and tense endgames stand out.
\end{enumerate}

\section{Limits}

The bands are a cross-section, not a path: the band above is other players, in games against
their own level, and their numbers say where stronger players give away less, not what he would gain
by fixing one thing. The rule's discount fit adjusted only for the rating difference, and the
results show what that left out (Section~\ref{sec:discount}); the fits with the opponent's points
came after the results were read and are descriptions, not tests. With no engine line, the reader
tags only a piece left hanging and a mate; a fork or a pin that wins material goes untagged. Lichess
evaluations come from its own server analysis, games someone asked to have analysed are not a random
sample, and his Lichess level is the app's conversion of his chess.com rating. His games come from
the app's reviews, with another engine and depth, so the paper sets no band's level beside his.
Players recur across games and no username is kept, so the intervals, drawn over games, are somewhat
narrow.

\end{document}
